\section{Adoption, Ownership, and Deployment Gates}
\label{sec:adoption}

KedgeFlow can be introduced around a selected protected action while existing models, graphs, and orchestration remain in place. OpenFlow, Kubernetes admission, and service-mesh control/data-plane separation offer architectural precedents for adding programmable enforcement to existing systems~\cite{openflow151,kubernetesadmission2026,istioarchitecture2026}. They do not predict KedgeFlow's market trajectory or make it a drop-in security boundary. Adoption depends on effective attachment points, credential and writer coverage, target capabilities, operational reliability, and measured integration cost.

\subsection{Buyers and Operational Owners}
A deployment needs both an organizational sponsor and explicit owners for its architectural roles. Table~\ref{tab:adoption-owners} connects three common interests to acceptance evidence. These are proposed adoption personas, not measured market segments.

\small
\begin{table}[!ht]
\begin{tabularx}{\linewidth}{@{}L{0.22\linewidth}L{0.31\linewidth}X@{}}
\toprule
\textbf{Sponsor / Owner} & \textbf{Reason to Adopt} & \textbf{Evidence Needed to Buy or Operate}\\
\midrule
Security and governance & Restrict confused-deputy authority paths; reconstruct who proposed, admitted and executed each effect. & Verified credential/native-path isolation, subject/workload delegation, trusted approval, retention controls and independently checkable lineage. NIST SP 800-207 can inform the design; this paper does not certify compliance.\\
Platform engineering and SRE & Coordinate competing changes and preserve recoverable operation state. & Predicate/writer coverage, atomic authority claims, target guard/completion scope, fault/restore tests and explicit unknown/compensation outcomes.\\
Application and agent developers & Reuse policy, admission, evidence and outcome services across existing workflows. & Working Agent Adapter integration, typed tool descriptors, supported remediation/retry semantics, bounded latency and usable failure handling. Shared services reduce duplicated work but do not eliminate resource-adapter implementation.\\
\bottomrule
\end{tabularx}
\caption{Proposed buying and operating interests, with implementation evidence for each. NIST zero-trust guidance supplies architectural principles, not an automatic product certification~\cite{nist800207}.}
\label{tab:adoption-owners}
\end{table}
\normalsize

Appoint one deployment owner accountable for acceptance and incident coordination. Security owns policy and approval trust; platform/SRE operates enforcement, authority state and recovery; the resource owner accepts policy correctness, predicate scope and writer coverage. Application teams own runtime integration. The product owner defines the promised task behavior, interruption/support experience and acceptable operating cost.

Named deputies and on-call escalation prevent stalled or unknown operations from falling between teams. Emergency administrative paths, often called \emph{break-glass} paths, have separately administered authority, explicit scope, auditing and coordination assumptions.

\subsection{Start with One Governed Action}
\label{sec:first-adoption-slice}

A first pilot should use a consequential task whose expected user value
justifies the engineering and interruption cost. Use one typed action, one supported runtime attachment,
one isolated target, and a target predicate and outcome that the team can verify.
Prefer an action whose native conditional update protects the required state.
If the policy has no wider shared invariant, a separate capacity reservation
and saga are unnecessary; exact admission, conserved grant issuance and atomic
claim remain required. Human review is required only by the chosen policy.
Hardware assurance is required when that policy demands it.

Nine logical roles do not require nine new services. For example, a team can
co-locate the five control roles, place the gateway and resource adapter in a
protected execution service, and reuse an existing evidence source. This reduces
deployment units, not authority checks. The harness must still lack equivalent
target privilege, and issuance and claim must share coordinated durable state.
Choose another action if those boundaries cannot be established.

\small
\begin{table}[!ht]
\begin{tabularx}{\linewidth}{@{}L{0.20\linewidth}L{0.35\linewidth}X@{}}
\toprule
\textbf{Work Package} & \textbf{Reuse or Reference Starting Point} & \textbf{Adopter Deliverable}\\
\midrule
Runtime integration & Existing documented hook or MCP client; supplied signed schemas/fixtures. & Agent Adapter for the selected invocation paths; exact proposal and stable operation identity; supported remediation/status handling.\\
Policy and trust & Existing IdP, policy engine, telemetry or review service; bounded Cedar/WebAuthn examples. & Trusted subject/workload enrollment, policy owner, provider verification, required review assurance and key lifecycle.\\
Authority and enforcement & Atomic storage and protected gateway; proxy claim/binding compatibility tests. & Production grant issuer, one issuance/claim domain, independent credential custody, target authorization and bypass denial.\\
Target and recovery & Existing conditional API and status endpoint; mock adapter and recovery tests illustrate the contract. & Typed Resource Adapter, declared predicate/completion level, authoritative outcome lookup, retention and on-call recovery.\\
\bottomrule
\end{tabularx}
\caption{A first-slice bill of materials. Reuse reduces new code; the right-hand
column remains accountable deployment work. The supplied service and proxy
examples have distinct profiles and do not form a turnkey production stack.}
\label{tab:adoption-work-packages}
\end{table}
\normalsize

Before pilot approval, the application owner states the protected scope and the
platform, security and resource owners accept their work packages. Record
workload-specific budgets for integration effort, admission delay, reviewer
wait and volume, service availability, unresolved-attempt escalation, and
operating/support cost. The pilot measures these budgets against the existing
workflow; exceeding them triggers scoped redesign or postpones expansion.
There is no universal numerical target in this paper.

Measure the baseline and pilot on the same task population. Report added
decision delay at the adopter's chosen percentiles separately from human-review
queue time and target completion time; record integration effort per runtime
and resource, review prompts per governed task, unresolved-attempt age and
support time. A denial count alone is not a quality metric: assess incorrectly
blocked authorized tasks against labeled policy cases. Expansion requires both
the boundary tests and acceptable user/operating cost.

\subsection{Adoption Friction and Its Response}
\label{sec:adoption-friction}

Table~\ref{tab:adoption-frictions} makes the cost of each boundary visible and
assigns a response. A mitigation may reduce engineering or user effort, but
cannot waive a contract required by the selected protected action.

{\small
\begin{longtable}{@{}L{0.23\linewidth}L{0.40\linewidth}L{0.32\linewidth}@{}}
\caption{Adopter friction, an owned mitigation, and evidence for expansion.}
\label{tab:adoption-frictions}\\
\toprule
\textbf{Friction / Owner} & \textbf{Practical Response} & \textbf{Acceptance Measure or Stop Condition}\\
\midrule
\endfirsthead
\toprule
\textbf{Friction / Owner} & \textbf{Practical Response} & \textbf{Acceptance Measure or Stop Condition}\\
\midrule
\endhead
\bottomrule
\endfoot
Too many components / platform lead & Assign logical roles to existing services before buying new ones; co-locate compatible roles and use one declared authority domain. Keep permission and key ownership explicit. & Named owners and a priced build/buy/retain work plan. Stop if required role ownership or support is missing.\\
\addlinespace
Harness changes and tool churn / application team & Start at one documented attachment; freeze typed inputs and reuse the same operation on remediation. Pin runtime/profile/descriptor versions and run the contract fixtures on upgrade. & Covered-path and substitution tests; measured integration effort. Unsupported paths stay disabled for the protected scope.\\
\addlinespace
Credential migration / security and platform teams & Provision enrollment, broker and protected service custody before cutover; revoke equivalent agent credentials only after the isolated route works. Inventory native tools, sockets and other writers. & Direct-bypass, revocation and secret-access tests. A working proxy without custody/coverage is still a pilot.\\
\addlinespace
Legacy target limitations / resource owner & Govern one typed API action; reuse a native conditional guard where it covers the predicate. Declare B/C exposure and recovery explicitly. & State-change and outcome-query tests. Block a protected action if its required predicate cannot be enforced.\\
\addlinespace
Approval fatigue / policy and product owners & Request review only for policy-required obligations; reuse still-valid static checks through non-authorizing leases. Route one canonical pending proposal to a reviewer and resume it asynchronously. & Review volume/wait, repeated prompts and authorized-task abandonment within adopter-set budgets; no blanket or stale approval reuse.\\
\addlinespace
Latency and dependency outages / SRE & Overlap permitted planning with admission; reuse authenticated broker connections and budget exchange plus final dispatch checks within signed validity. Measure the complete route, including optional grant-store lookup and header limits. & Admission/exchange tail latency, evidence age, remaining dispatch margin and fault tests. Expiry never permits lifetime extension or fresh authority. Required admission/claim failure blocks new protected dispatch.\\
\addlinespace
Ambiguous failures and support load / SRE and resource owner & Retain durable operation/attempt IDs, correlate authoritative target status, expose pending/unknown state and set bounded polling with an owned escalation deadline. Qualify recovery before admitting legacy actions. & Lost-response, late-effect and escalation tests; unresolved age and operator time. Deadline or incident closure cannot clear consumption or justify an unsafe inverse.\\
\addlinespace
Upgrade cost and lock-in / platform and procurement owners & Require versioned descriptors, exportable contracts/lineage and an explicit key/state transition plan. Qualify replacements against the same scope and fault cases. & Upgrade and migration rehearsal; no transparent conversion of signed records, replay of consumed roots or untested portability claim.\\
\end{longtable}
}


\subsection{A Four-Stage Rollout with Explicit Gates}
Figure~\ref{fig:adoption-stages} outlines a candidate rollout. The stages describe growing verified coverage, not permission to weaken the execution contract on an already protected action. Single-use claim is required as soon as grants authorize execution; shared-invariant reservations are added when the declared predicate requires them.

\begin{figure}[!ht]
\centering
\begin{tikzpicture}[>=Latex,node distance=0.55cm,
 box/.style={draw=acpline,fill=acppale,rounded corners=4pt,align=left,text width=13.0cm,inner sep=7pt,font=\small},
 gate/.style={font=\scriptsize\sffamily,fill=white,align=center,inner sep=2pt}]
\node[box] (observe) {\textbf{1. Observe: no enforcement claim}\\
 Inventory selected calls and candidate policy decisions; record capture gaps and uncertainty.\\
 \emph{Gate:} bounded telemetry, protected retention, known invocation coverage.};
\node[box,below=of observe,draw=acpblue] (pilot) {\textbf{2. Pilot: enforce a selected isolated route}\\
 Test typed proposals, authenticated remediation, admission, grants, atomic claim and outcomes.\\
 \emph{Gate:} valid/invalid flows pass; unsupported clients and paths cannot bypass the pilot target.};
\node[box,below=of pilot,draw=acpamber] (cutover) {\textbf{3. Cut over privilege: protect a declared production surface}\\
 Remove equivalent agent target authority; enable independently protected enforcement and recovery.\\
 \emph{Gate:} deployed OS/network/target coverage, identity trust, fault handling and resource guarantees.};
\node[box,below=of cutover,draw=acpteal] (coordinate) {\textbf{4. Expand coordination: more writers, resources and workflows}\\
 Add full-scope reservations, partitioned saga allowance and supported compensation where needed.\\
 \emph{Gate:} competing-writer, conservation, late-effect, recovery and restore tests for each added scope.};
\draw[->,thick,acpgray] (observe)--(pilot);
\draw[->,thick,acpgray] (pilot)--(cutover);
\draw[->,thick,acpgray] (cutover)--(coordinate);
\end{tikzpicture}
\caption{A proposed rollout, with the assurance boundary stated at every stage. Shadow observation is not protected execution, and expanding a deployment cannot postpone the claim/credential rules of an existing protected route.}
\label{fig:adoption-stages}
\end{figure}

\paragraph{1. Observe.}
An Agent Adapter or existing tracing middleware records selected invocations and evaluates candidate policies without changing execution authority. Captured arguments can support a provisional proposal; missing identity, evidence or outcome data must remain missing. A log does not create a verified witness or target commit receipt. Observation can expose credentials or add overhead, so redaction, access controls, retention and performance are part of the pilot. This stage establishes an inventory and policy hypotheses, with no claim of prevented mutations.

\paragraph{2. Pilot a governed route.}
Use an isolated or nonproduction target and an explicitly supported runtime adapter. Exercise the existing \texttt{NEEDS\_EVIDENCE} contract, denied/unknown results, exact parameter and identity binding, grant issuance/claim, and the adapter's declared outcome query. The pilot must restrict its target authority sufficiently to test enforcement. A proxy on one call path can validate that route, but remaining native credentials elsewhere prevent a wider protection claim. Authorized observations use their own confidentiality and scope rules.

\paragraph{3. Cut over protected privilege.}
Revoke or restrict equivalent agent credentials for the chosen production resource/action surface; provision independently protected gateway/broker custody and target authorization. Verify alternate shell, file, socket, SDK and delegated paths, user/workload trust, policy/evidence freshness, credential downscoping and outage behavior. Preserve separately administered break-glass procedures. A cutover is accepted only when bypass and resource-specific fault tests pass; fallback to an ambient credential would undo the boundary.

\paragraph{4. Expand coordination and recovery.}
Add new resources and participating writers to declared conflict domains, and coordinate any shared invariant before expanding autonomous operation. Two grants in the same conflict domain need not be mutually exclusive if the invariant permits both; competing removals of the last spare capacity must be. Saga orchestration allocates conserved child allowances and runs registered, freshly admitted compensations only when known outcomes and current guards justify them. It is conditional recovery, not automatic atomic rollback. Unknown forward effects first require reconciliation and a late-effect barrier.

\subsection{Deployment Packaging and Guarantee Limits}
The architecture permits several packaging choices. A local developer distribution can combine the controller roles, a protected gateway, descriptors and compatibility vectors; single-host storage still needs the same atomic issuance/claim and recovery semantics. An enterprise service can centralize policy administration, review trust and retained lineage while gateways remain close to resources. A provider integration can supply workload enrollment, target-specific credentials and native guards. These are design options, not announced products, tested marketplace installers or revenue forecasts. Multi-region availability additionally needs explicit authority ownership, quorum/partition and restore contracts; deployment geography alone does not supply them.

CTDs let teams govern a selected consequential surface while retaining permitted workspace operations. The friction budgets in Section~\ref{sec:first-adoption-slice} are measured on the actual workload; neither locality nor static leases establish a cost or latency advantage. Leases reuse still-valid admission work but do not authorize execution or replace final freshness checks.

Legacy resources can use declared Level B or C contracts when the selected policy accepts their check-to-effect exposure or delayed uncertainty. If the required protected-action predicate cannot be enforced, the action must block until an adequate target mechanism exists; unenforced policy evaluation belongs only to the explicitly labeled observation stage. Protocol neutrality requires independently working adapters and common fixtures; MCP/A2A transport support alone is insufficient.

\subsection{Workflow Interruptions and Safe Next Actions}
\label{sec:adoption-workflow}

The Agent Adapter and product UI expose the existing result and lifecycle
semantics faithfully. Table~\ref{tab:adoption-handling} is integration guidance,
not a new wire error or state profile. Each protected task shows the exact
action, a policy reason safe to disclose, its proposal/attempt reference,
current status, responsible reviewer/operator and next supported action.
Deduplicate review prompts by stable proposal/operation identity. Editing a
reviewed action creates a new proposal and admission; a terminal confirmation
cannot substitute for a required hardware-review obligation.

{\small
\begin{longtable}{@{}L{0.22\linewidth}L{0.42\linewidth}L{0.31\linewidth}@{}}
\caption{Existing results mapped to workflow handling and user communication.}
\label{tab:adoption-handling}\\
\toprule
\textbf{Result or Lifecycle State} & \textbf{Workflow Response} & \textbf{User Status / Accountable Owner}\\
\midrule
\endfirsthead
\toprule
\textbf{Result or Lifecycle State} & \textbf{Workflow Response} & \textbf{User Status / Accountable Owner}\\
\midrule
\endhead
\bottomrule
\endfoot
\nolinkurl{NEEDS_EVIDENCE} & Pause the unchanged proposal; validate the trusted remediation path and obtain only the required evidence. Resume exact admission rather than replanning silently. & Awaiting named evidence/review, with expiry and a status reference. Application/product owner and review provider.\\
\addlinespace
\nolinkurl{DENIED}; identity/grant mismatch & Stop protected dispatch. Route policy changes through authorized administration; investigate identity/binding errors. Do not repeatedly ask another reviewer for the same refused action. & Denied or binding error, safe reason and escalation path. Policy/security owner.\\
\addlinespace
Stale policy/evidence; \nolinkurl{STATE_GUARD_VIOLATION} & Refresh and revalidate through admission. Inspect existing root/attempt state before issuing replacement authority; changed bound context or parameters require the appropriate new proposal. & Action needs revalidation, not completed. Admission/resource owner.\\
\addlinespace
\nolinkurl{ROOT_UNAVAILABLE}; \nolinkurl{NONCE_CONSUMED} & Query the existing authority/attempt. A new RPC ID, nonce or automatic tool retry cannot create an additional allowance. & Existing operation or authority issue; status lookup and operator escalation. Platform owner.\\
\addlinespace
\texttt{dispatched} / \texttt{pending} & Query the correlated target operation under its completion contract; do not redispatch to obtain a faster response. & Accepted or executing, with no premature completion claim. Resource adapter/SRE.\\
\addlinespace
\texttt{unknown}; \nolinkurl{OUTCOME_UNKNOWN} & Retain consumption and any reservation; reconcile authoritative evidence and escalate at the declared deadline. Another action does not clear uncertainty. & Outcome unknown, last confirmed fact and incident reference. Resource owner/SRE.\\
\addlinespace
\texttt{committed}; \texttt{failed} / \texttt{compensation\_failed} & Report the receipt's actual effect scope. Failure does not imply zero effect or refund; any supported compensation needs fresh admission and known preconditions. & Confirmed effect or failure with possible partial effect. Resource/product owner.\\
\addlinespace
\texttt{not\_dispatched}; \texttt{fenced\_no\_effect} & Preserve the verified refusal/resolution; \texttt{not\_dispatched} alone does not restore an allowance. Authoritative no-effect evidence plus the required late-effect barrier can support restoration under issuer conservation rules. & No dispatch or fenced no effect, with supporting evidence. Gateway/resource owner.\\
\end{longtable}
}

Before claim, UI cancellation may withdraw a proposal or request atomic
retirement of unclaimed authority under the authority contract. For claimed
or dispatched work, cancellation stops future work only under the declared
target contract: it does not prove no effect, release a reservation or undo a
commit. Without authoritative resolution and any required late-effect barrier,
the attempt remains pending or unknown. The product owner explains this
distinction before launch.

\subsection{Operate, Upgrade, and Exit Without Reopening Authority}
\label{sec:adoption-operations}

During an outage, permitted planning and workspace work can continue within
their policy. New protected dispatch stops when required admission, claim or
credential authority is unavailable. An Outcome Service outage may buffer
already supported records in durable outboxes under the declared contract;
if a required transition record cannot be durably retained, that transition
must stop. Disconnected operation requires a supported local authority/target
profile preserving custody, freshness, atomic claim and recovery. Cached policy
or admission alone is insufficient.

The deployment owner rehearses cutover, abort, upgrade, restore and supplier
exit as one authority-preserving transition:
\begin{enumerate}
\item Quiesce new issuance and dispatch; establish the resource's fencing or
late-effect barrier where required. Inventory every granted, claimed, pending
and unknown attempt. Stopping the controller does not cancel target work.
\item Reconcile or quarantine in-flight operations. Retain consumed roots,
generations, reservations, replay state, dispatch intents and outboxes; a stale
restore or a fresh empty ledger cannot make old authority reusable.
\item Preserve original signed records, trusted key/version histories and
descriptor/policy bindings. Validate the new runtime/profile and key transition;
incompatible signed records are not silently rewritten or re-signed as proof
of the old history.
\item Run the same scope-specific refusal, bypass, concurrency, lost-response
and restore cases before resuming. Separately authorized emergency operators
remain declared writers. Returning equivalent ambient credentials to an agent
changes or ends the protected scope and requires explicit owner acceptance.
\end{enumerate}

The acceptance record names dependency patching/provenance, key rotation,
outage backlog, recovery deadlines and on-call ownership. A failed pilot can
disable the protected feature while retaining permitted work and a separately
governed operator procedure. This makes postponement practical without treating
an unsafe fallback as a successful migration.

\subsection{Buyer Control and Supplier Evidence Record}
\label{sec:buyer-control-record}

All buyers need a deployment-specific evidence record. Government and
sovereign procurement teams additionally use it to state who controls the
deployment and its dependencies; local hosting alone does not establish that
control. Name the operators and administrators of policy, grant signing,
workload enrollment, Evidence Providers, credential brokers, optional grant
stores/resolvers, gateways,
authority state and retained outcomes. Record which external model, judge,
identity, signing or support service can observe, alter or disable the chosen
operation, and whether the buyer can operate the required path independently.
Declare provider independence and conflicts of interest rather than inferring
them from multiple signatures.

Inventory the sensitive canonical proposals, identity, evidence, review and
receipt records: access, storage/transfer location, retention, permitted
disclosure and export. Redacted views retain links to the protected originals;
altering a signed payload breaks its signature and digest bindings. Specify
who may verify the retained chain, what records verification requires, and how
originals/key histories survive supplier exit.

\small
\begin{table}[!ht]
\begin{tabularx}{\linewidth}{@{}L{0.22\linewidth}L{0.34\linewidth}X@{}}
\toprule
\textbf{Acceptance Topic} & \textbf{Evidence Supplied Here} & \textbf{Evidence Required from Deployment}\\
\midrule
Contract exchange & Closed schemas, signed vectors and independent canonicalization/signature checks. & Frozen versions and independently implemented peers exchanging the selected profile.\\
Protected boundary and review & Proxy identity/claim tests; mock guards; software-key WebAuthn mechanics. & Actual enrollment, hardware assurance when required, native bypass/custody and selected-target guard/outcome tests.\\
Recovery and operation & Bounded claim races and mock lost-response reconciliation; architecture acceptance cases. & Real fault/restore, key transition, emergency operation, data controls, support and safe supplier-exit rehearsal.\\
Adoption value & Proposed work packages, friction responses and measurement method. & Workload-specific user/operating cost, accepted guarantee limits and signed-off expansion decision.\\
\bottomrule
\end{tabularx}
\caption{Procurement evidence split. Section~\ref{sec:executed-checks} states the
scope of supplied checks; Section~\ref{sec:acceptance} supplies deployment tests.
This is an acceptance input, not a certification or jurisdiction-specific
accreditation claim.}
\label{tab:buyer-evidence}
\end{table}
\normalsize

The resource/public-service owner accepts policy correctness and the residual
limits of B/C adapters, external writers and uncertain effects. The buyer
selects any legal, national-assurance or accreditation requirements separately
and obtains evidence for its chosen profile. No sovereign-operation,
production-resilience or supplier-portability result is established by this
whitepaper.


\subsection{Application-Team Onboarding Record}
Before accepting a protected deployment, retain one reviewable record containing:
\begin{enumerate}
\item Resource/action scope, CTDs, all relevant writers and credential/socket/native-tool routes, including excluded or separately administered paths.
\item One accountable deployment owner and named component/interface, product-support and incident owners; trust enrollment, secret-store access policy, active runtime/profile/schema/descriptor versions, policy epoch, gateway audience, resource incarnation and authority domain.
\item Required obligation classes, review assurance, authority/claim domain, target predicates, completion modes and compensation limits.
\item Executed missing-grant, forged/expired evidence, identity/parameter substitution and direct-bypass cases, with no target dispatch on refusal.
\item Competing-writer, concurrent-redemption, lost-response, late-effect and restore tests; consumed unknown authority cannot reopen through a new request ID.
\item Measured operating/user cost, data/dependency control and the rehearsed outage/upgrade/exit procedure, with an explicit list of partial or unverified contracts. Section~\ref{sec:acceptance} supplies the technical cases and Section~\ref{sec:buyer-control-record} the buyer evidence split.
\end{enumerate}

This record gives buyers a concrete completeness assessment and gives implementers an integration backlog. Each new tool is assigned by the obligation, owned authority and failure contract it supplies, preserving the component vocabulary as the ecosystem changes.
